\section{The Parallel Engine}\label{sec:engine}

This section describes the engine as implemented, because the diagnosis in
\cref{sec:diagnosis} is a statement about \emph{this} design.  Nothing here is
claimed as novel; it is the design a careful implementer arrives at, and the
point of the paper is that arriving at it is not sufficient.

\subsection{Purity and level-synchronous exploration}

The parallelisation rests on one observation.  $\succfn(q,Z)$ is a \emph{pure
function}: its inputs are the read-only network and the pair $(q,Z)$, its output
is a freshly allocated vector of successors, and it neither reads nor writes
$\Passed$ or any other state's data.  Calls to $\succfn$ on distinct states are
therefore independent.

What is \emph{not} independent is the check-and-insert that follows each call.
\Cref{alg:seq} interleaves the two per iteration, and that interleaving is what
forces sequential execution.  Separating them by BFS level breaks the
dependency: states are grouped into a frontier $\Frontier_k$, and every state in
$\Frontier_k$ was already checked against $\Passed$ before being placed there, so
their expansions may run in any order (\cref{fig:par-bfs}).

\subsection{A lock-striped shared passed list}

A natural next step---and the one our earlier design took---is to have each
thread expand into a private buffer and then merge all buffers sequentially.  We
abandoned it.  A sequential merge reintroduces a serial phase proportional to the
number of \emph{successors} (roughly $10\times$ the number of states), and it
forces a second pass over data that has just left cache.

Instead, expansion and insertion are \emph{fused}: a thread inserts each
successor into the shared passed list as soon as it is generated.  The passed
list is sharded into $S = \max(1024, 64T)$ \emph{stripes} by a multiplicative
hash of the discrete key $q$, each stripe carrying its own mutex.  A thread
therefore holds one stripe lock for the duration of one antichain update, and
$S \gg T$ keeps collisions rare.

Per stripe the table is flat and open-addressed: a slot array maps $q$ to the
head of its antichain, and antichain nodes live in a deque with stable addresses,
each holding an index into the shared zone-intern table plus the BFS level that
created it.  Storing the discrete part once per bucket (rather than per zone) and
referring to zones by index keeps the resident frontier small enough to stay
cache-resident: a frontier entry is $12$~bytes, not a full DBM.

\begin{algorithm}[t]
\caption{Level-synchronous parallel BFS with a striped shared passed list}
\label{alg:par}
\begin{algorithmic}[1]
\Require network $\mathcal{N}$, target $\varphi$, thread count $T$
\State $\Passed[q_0] \leftarrow \{Z_0\}$;\enspace
       $\Frontier \leftarrow [(q_0, Z_0)]$;\enspace $k \leftarrow 0$
\While{$\Frontier \neq \emptyset$}
  \State $k \leftarrow k+1$
  \ParFor{$i = 0,\ldots,|\Frontier|-1$}
    \Comment{\texttt{schedule(dynamic,8)}}
    \ForAll{$(q', Z') \in \succfn(\Frontier[i])$}
      \If{$\varphi(q', Z')$} \Return \textbf{true} \EndIf
      \State \textbf{lock} stripe of $q'$;\enspace
             $\mathrm{add}(q', Z', k)$;\enspace \textbf{unlock}
             \Comment{antichain update, \cref{alg:add}}
    \EndFor
  \EndParFor
  \State $\Frontier \leftarrow$ entries of $\Passed$ stamped with level $k$
         \Comment{gathered in parallel}
  \State reclaim retired antichain nodes and unreferenced interned zones
\EndWhile
\Return \textbf{false}
\end{algorithmic}
\end{algorithm}

\subsection{The antichain update}

\Cref{alg:add} is where the time goes, so we state it precisely.  Given an
incoming $(q', Z')$ it performs three folds over $q'$'s antichain---drop $Z'$ if
some stored zone subsumes it; retire stored zones that $Z'$ subsumes; absorb
stored zones whose union with $Z'$ is exactly convex---and finally interns and
prepends $Z'$.

\begin{algorithm}[t]
\caption{$\mathrm{add}(q, Z, k)$ --- antichain update (original)}\label{alg:add}
\begin{algorithmic}[1]
\If{$q$ is new} intern $Z$, create its antichain, \Return \EndIf
\ForAll{live nodes $m$ of $q$'s antichain}
  \Comment{(a) subsumed? -- \emph{one unpack per node}}
  \If{$Z \preceq_{aLU} \mathrm{unpack}(m)$} \Return \textsc{dropped} \EndIf
\EndFor
\ForAll{live nodes $m$}
  \Comment{(b) retire nodes subsumed by $Z$}
  \If{$\mathrm{unpack}(m) \preceq_{aLU} Z$} mark $m$ retired \EndIf
\EndFor
\While{some live $m$ has $Z \cup \mathrm{unpack}(m)$ convex}
  \Comment{(c) convex-union absorption}
  \State $Z \leftarrow Z \sqcup \mathrm{unpack}(m)$; mark $m$ retired
\EndWhile
\State intern $Z$ and prepend it to the antichain
\end{algorithmic}
\end{algorithm}

Every $\mathrm{unpack}$ in \cref{alg:add} is an $O(n^3)$ closure
(\cref{sec:mingraph}), and fold~(c) restarts its scan after each absorption, so a
node can be unpacked repeatedly.  \Cref{sec:diagnosis} measures the consequence.

\subsection{Between levels}

After the parallel region, survivors are gathered into the next frontier in
parallel over the touched discrete keys.  Two reclamation passes then run:
retired antichain nodes are compacted away per stripe (one thread per stripe,
geometrically triggered), and the zone-intern table is mark-and-swept so that
records no longer referenced by any live node or frontier entry return to a free
list.  The second of these is where \cref{sec:gc} finds its Amdahl term.
